\section{Indicators of Compromise (IOCs)}
\label{sec:ioc}

The following indicators allow detection of this compromise on the live host and across other
instances of the application.

\subsection{On-Disk Indicators}
\begin{itemize}
  \item \textbf{Known artifacts:} \fillme{path/to/artifact} and any file matching
        \fillme{describe the general pattern, e.g. "an unexpected .jsp/.php under an
        upload/attachment directory"}.
  \item \textbf{Filename tokens:} \fillme{list suspicious tokens, e.g. dot-named files, double
        extensions, null bytes, oddly-cased extensions}.
  \item \textbf{Content markers:} \fillme{list suspicious code markers, e.g. Runtime.exec,
        ProcessBuilder, base64 decode-and-eval patterns}.
\end{itemize}

\begin{codebox}[Read-only on-disk detection sweep]
# illustrative grep/find command for hunting the same pattern elsewhere -- replace with your own
\end{codebox}

\subsection{Network / HTTP Indicators}
\fillme{Note whether access logs were available. If not available in this source tree, say so and
list what to pull and search for.}
\begin{itemize}
  \item \fillme{e.g. GET requests to the known artifact path, especially with suspicious parameters}
  \item \fillme{e.g. POST/multipart uploads with executable-looking filenames}
  \item \fillme{e.g. outbound connections initiated by the application process}
\end{itemize}

\newpage
\subsection{Host Indicators}
\begin{itemize}
  \item \fillme{e.g. unexpected database logons sourced from the application host}
  \item \fillme{e.g. modifications to authorized\_keys, cron, or service startup scripts}
  \item \fillme{e.g. new or modified files anywhere under the web root, especially world-writable ones}
\end{itemize}
